\begin{document}

\subsection{Pástor & Veronesi (2009) — "Technological Revolutions and Stock Prices"}

The central question of this paper is: why do stock prices surge when a new technology arrives, and then often crash — even when the technology turns out to be genuinely valuable?
The conventional answer is "irrational bubble." Pástor and Veronesi argue something more subtle: this pattern can be entirely rational. Their model works through three key mechanisms.
First, when a major new technology appears (electricity, the internet, AI), nobody knows whether it will be widely adopted or how productive it will actually be. This uncertainty itself is valuable to stockholders — it creates an option-like payoff. Stock prices rise to reflect not just expected productivity but the sheer range of possible outcomes, because the upside scenarios are enormous.
Second, as time passes and more information arrives, uncertainty resolves. Even if the news is mostly good — the technology does get adopted — the reduction in uncertainty removes that option value. Stock prices can fall even as the technology succeeds. This is counterintuitive but powerful: you don't need bad news to get a crash, you just need the uncertainty to go away.
Third, the old economy matters too. When a new technology arrives, it's unclear whether incumbent firms will adopt it successfully or be displaced. This creates a discount rate channel — the riskiness of the entire economy shifts during the transition, which affects how future cash flows are valued.
The paper formalizes all of this in a general equilibrium model and calibrates it to match patterns from the late 1990s internet boom. The key empirical predictions are that stock prices rise when new technologies arrive, the initial rise is larger when uncertainty is greater, and prices eventually fall as uncertainty resolves — producing a "bubble-like" pattern from purely rational behavior.
For your thesis, this is the paper that keeps you honest. When you find explosive dynamics in AI stock prices using the PSY test, Pástor-Veronesi forces you to ask: is this irrational speculation, or is it rational pricing of genuine uncertainty about a transformative technology? That distinction is at the heart of what makes your project interesting.


\section{Scheinkman & Xiong (2003) — "Overconfidence and Speculative Bubbles"}
This paper asks: why would a rational investor knowingly pay more than fundamental value for an asset? The answer is the resale option.
The setup has three ingredients. First, agents are overconfident — they overweight the precision of signals they personally observe relative to signals others observe. This generates persistent disagreement about the asset's fundamental value, even though everyone sees the same public information. Second, there are short-selling constraints — pessimists can't easily bet against the asset, so only the optimists' views are reflected in the price. Third, and this is the key insight, agents are aware that other agents are also overconfident. So even if you think the asset is somewhat overvalued right now, you're willing to buy it because you expect that in the future, someone even more optimistic than you will come along and buy it from you at a higher price.
This creates a speculative premium on top of the fundamental value — the "bubble component." The asset price equals fundamental value plus a resale option value. That option is valuable because belief heterogeneity means there will always be a future state where someone disagrees with you enough to pay more. The bubble is larger when: overconfidence is greater, trading costs are lower, and uncertainty about fundamentals is higher.
The model also predicts excessive trading volume during bubble episodes — because agents keep trading back and forth as their relative optimism shifts. This is empirically testable and has been observed in AI stocks.
For your thesis specifically:
This paper gives you the behavioral explanation for why bubbles form, as opposed to Pástor & Veronesi's rational explanation. If your PSY tests detect explosive dynamics in high-AI-dependency stocks, Scheinkman-Xiong says: that's because investors disagree massively about AI's future payoff, pessimists can't easily short these stocks (especially Nvidia, where short interest has been relatively low), and each buyer expects to sell to a greater optimist later. The disagreement about AI's economic impact is enormous right now — some analysts project transformative productivity gains, others see it as overhyped infrastructure spending — so the conditions for the Scheinkman-Xiong mechanism are clearly present.
You can also use their trading volume prediction. If you find that bubble episodes identified by PSY coincide with abnormally high trading volume in your AI-exposed stocks, that's consistent with the overconfidence/resale-option story and harder to explain through Pástor & Veronesi's rational framework alone.

\end{document}
